\begin{textsample}{POS Dim 3 – persona\_gpt – Score 32.00 – t596\_gpt.txt}  \label{ex:f3_pos_024}
Europe is running out of \textbf{time} to tackle the climate crisis, and the way we move is at the heart of the \textbf{problem}.
The \textbf{transport} \textbf{sector} is responsible for more than a quarter of Europe ’s CO2 \textbf{emissions}, and almost \textbf{half} of that – 45 % – comes from road \textbf{transport} alone.
Unless we change how we travel, we \textbf{will} miss the narrow window we have in the next decade to keep global warming in check.
Phasing out diesel and petrol vehicles is central to this transformation.
Research by the Ecologic Institute and Greenpeace Germany lays out six key policy \textbf{measures} that can drive this \textbf{shift} and help rapidly cut \textbf{emissions} from \textbf{transport}.
First, governments \textbf{need} to set a clear \textbf{end} date for the \textbf{sale} of new internal combustion engine vehicles and back it up with binding \textbf{production} quotas for electric vehicles ( EVs ).
This provides certainty for \textbf{industry}, accelerates \textbf{innovation}, and ensures that cleaner vehicles become the norm rather than the exception.
Second, \textbf{public} authorities must rapidly expand and subsidize charging \textbf{infrastructure}.
People \textbf{will} only switch to EVs if charging is easy, affordable, and widely available – in cities, along highways, and in rural areas. \textbf{Public} \textbf{investment} and smart subsidies can close the gaps and make EV charging as straightforward as filling up a tank.
Third, incentives \textbf{need} to evolve with the \textbf{market}.
Early on, generous support can help kick-start EV adoption ; later, as EVs become more common and cheaper, incentives should be adjusted and better targeted.
This avoids \textbf{wasting} \textbf{public} \textbf{money} and keeps support focused where it ’s most effective – for example, helping low- and middle-income \textbf{households} make the switch.
Fourth, no single policy is enough.
A mix of tools – such as subsidies, mandates, \textbf{standards}, and \textbf{regulations} – works \textbf{better} than any one \textbf{measure} on its own.
Combining carrots and sticks helps overcome different barriers at the same \textbf{time}, from upfront \textbf{costs} to \textbf{supply} constraints.
Fifth, governments should make polluters pay while rewarding cleaner choices.
That \textbf{means} higher \textbf{taxes} or charges on fossil-fuel vehicles and fuels, alongside financial advantages for EVs.
Well-designed \textbf{taxes} and incentives can \textbf{shift} both \textbf{consumer} demand and manufacturer decisions, speeding up the transition away from diesel and petrol.
Finally, clear communication is essential.
People \textbf{need} to understand not just what is changing, but why it matters and how it \textbf{will} improve their lives – from cleaner air and quieter streets to lower running \textbf{costs} and more stable \textbf{energy} \textbf{prices}.
Transparent, consistent messaging helps build trust and \textbf{public} support for ambitious climate action in \textbf{transport}.
Electric vehicles are a crucial part of the \textbf{solution}, but they are not the whole story.
The goal is not simply to replace every diesel and petrol \textbf{car} with an electric one.
To truly cut \textbf{emissions} and build livable cities, Europe must also reduce the overall \textbf{number} of private \textbf{cars} on the road and prioritize \textbf{alternatives} : high-quality, affordable \textbf{public} \textbf{transport} powered by renewable \textbf{energy}, as well as safe and attractive \textbf{infrastructure} for \textbf{cycling} and walking.
The next decade is decisive.
By combining a rapid phase-out of fossil-fuel \textbf{cars} with strong support for EVs and a broader \textbf{shift} towards \textbf{public} \textbf{transport}, cycling, and walking, Europe can slash \textbf{transport} \textbf{emissions} and move decisively toward a cleaner, fairer future.

% matched lemmas: alternative, car, consumer, cost, cycle, emission, end, energy, good, half, household, industry, infrastructure, innovation, investment, market, mean, measure, money, need, number, price, problem, production, public, regulation, sale, sector, shift, solution, standard, supply, tax, time, transport, waste, will
\end{textsample}
